\chapter{Literature Review}
\label{ch:literature}

\section{Oil Prices and Macroeconomic Performance}

The modern literature on oil shocks and macroeconomics begins with Hamilton (1983), who showed that most US recessions since World War II were preceded by increases in the price of crude oil. That finding prompted a large body of work examining the channels through which oil prices affect output, employment, and inflation. Early studies treated oil shocks as straightforward supply-side events: higher energy costs raise production costs, squeeze profit margins, and push up the general price level. This supply-side mechanism is intuitive and remains a standard feature of macroeconomic textbooks.

By the mid-1990s, a more nuanced picture had emerged. Mork (1989) split oil price changes into positive and negative movements and found that oil price increases depressed US GDP growth, but oil price decreases did not produce a symmetric expansion. Lee, Ni and Ratti (1995) reinforced this finding by showing that oil price volatility conditions mattered: a price increase that comes during a period of stability has more macroeconomic impact than the same increase during a volatile period. Hamilton (2003) formalised various nonlinear measures of oil shocks and argued that the net oil price increase specification better captures the recessionary effects of oil in the US. These studies established asymmetry as a central theme in oil shock research.

On the inflation side, the initial empirical consensus was clear: oil price increases raised inflation. But Hooker (2002) challenged this consensus by showing that the oil--inflation relationship weakened after 1980 in the United States. He attributed this to improved monetary policy credibility, which anchored inflation expectations and prevented oil price increases from generating persistent inflationary spirals. Blanchard and Gal\'{i} (2009) provided additional evidence across developed economies, documenting a decline in the pass-through of oil shocks to both inflation and output since the mid-1980s.

\section{Asymmetric Pass-Through: Rockets and Feathers}

The idea that prices respond faster to cost increases than cost decreases originates in Bacon's (1991) study of UK petrol markets, which gave rise to the ``rockets and feathers'' label. The intuition is straightforward: firms raise output prices quickly when input costs rise, but lower them slowly, or not at all, when costs fall. The pattern has been documented in retail fuel markets across many countries, though the evidence is not unanimous.

At the macroeconomic level, asymmetry in the oil--inflation relationship has been studied using various econometric approaches. Lardic and Mignon (2008) used asymmetric cointegration techniques to study oil price effects in G7 economies and found evidence that the long-run relationship between oil prices and GDP differed depending on the sign of the oil price change. Shin, Yu and Greenwood-Nimmo (2014) developed the nonlinear autoregressive distributed lag (NARDL) framework, which decomposes a regressor into positive and negative partial sum components within a cointegrating equation. Their approach has become popular in applied work because it provides a test for long-run asymmetry within a single-equation error correction framework.

However, not all studies support the asymmetry hypothesis. Kilian and Vigfusson (2011) argued that much of the evidence for asymmetric macroeconomic effects of oil shocks is based on censored regressions that introduce a pre-test bias. Using properly specified models and US data, they found that the evidence for asymmetric effects on GDP was weaker than commonly believed. This critique is relevant for inflation pass-through studies as well, because the same statistical issues apply when oil price changes are split into positive and negative components.

\section{Oil Pass-Through in Emerging Economies}

The pass-through literature is heavily concentrated on developed economies, particularly the United States and Europe. Emerging economies, however, face different conditions that may amplify or dampen oil pass-through. Higher energy intensity of GDP, larger shares of food in consumption baskets, more volatile exchange rates, and fuel price regulation all affect how oil shocks transmit to consumer prices.

Chen (2009) studied oil price pass-through to inflation across 19 industrialised countries and found considerable cross-country variation. Countries with less credible monetary policy and more rigid labour markets tended to exhibit stronger pass-through. For developing countries, the evidence is thinner but growing. Several studies have used vector autoregression and structural approaches to estimate oil price effects, but results vary considerably depending on the time period, the oil price measure, and the identification strategy.

\section{India-Specific Studies}

India's case is particularly interesting because of its high oil import dependence and its history of administered fuel pricing. Bhanumurthy, Das and Bose (2012) examined the macroeconomic effects of oil price shocks on India using a macroeconometric model and found that oil price increases raised inflation and widened the current account deficit. However, they noted that the government's pricing policy at the time, which involved substantial fuel subsidies, significantly dampened the pass-through.

Abu-Bakar and Masih (2018) applied the NARDL framework to India and found evidence of asymmetric oil price pass-through to inflation. Their study covered a period when fuel pricing was partly regulated, and they argued that the asymmetry arose because the government was quicker to raise retail prices when crude prices went up than to lower them when crude prices fell. This is a plausible mechanism, though it is harder to document precisely at the monthly frequency.

Pradeep (2022) studied the effect of diesel price deregulation on the asymmetry of oil price pass-through in India. Using a nonlinear ARDL approach, Pradeep found that deregulation made the pass-through more symmetric in the post-reform period, because retail prices began tracking international prices in both directions. This finding has direct relevance for the present study, because the sample period here covers both the pre-deregulation and post-deregulation regimes.

The Reserve Bank of India itself has examined the oil--inflation link in its Monetary Policy Reports. The RBI has consistently noted that crude oil prices are a key source of inflation risk, particularly through their effect on transport costs and the prices of oil-intensive manufactured goods. The RBI's own estimates suggest that a 10 per cent increase in crude oil prices raises headline CPI inflation by about 0.3--0.4 percentage points over a one-year horizon, though this estimate includes both direct and indirect effects.

\section{The Exchange Rate Channel}

A feature of oil pass-through in India that deserves separate attention is the exchange rate. Since India buys oil in US dollars, the rupee--dollar exchange rate mediates the domestic-currency cost of imported crude. A rise in global oil prices combined with a weakening rupee compounds the inflationary impact. Conversely, a fall in oil prices may not fully reduce domestic fuel costs if the rupee depreciates simultaneously.

This dual channel, oil prices in dollars and the exchange rate, means that studies using a composite rupee oil price may conflate two distinct sources of inflationary pressure. If the rupee depreciates for reasons unrelated to oil (such as capital outflows or a widening current account deficit), the composite oil price will show a ``shock'' that is really an exchange rate event. Separating the Brent price in dollars from the exchange rate, as the robustness analysis in this dissertation does, provides a cleaner identification of each channel.

\section{Gaps in the Literature}

Several gaps motivate this dissertation. First, most India studies use data that stop before the full effect of diesel deregulation can be observed. This study extends the sample to December 2024, providing a long post-deregulation window. Second, many studies apply the NARDL or SVAR framework without checking whether simpler reduced-form models deliver similar conclusions. This study uses a transparent ADL framework in first differences, which is easier to interpret and less sensitive to cointegration rank assumptions. Third, few India studies explicitly separate the oil price channel from the exchange rate channel with both reported in the same specification. The Brent-plus-exchange-rate robustness model in this study addresses that gap. Fourth, using the sub-index for Fuel and Light CPI from the official MoSPI data provides a direct check on whether the dilution effect in headline CPI obscures the true pass-through.
